% ============================================================
%  CHAPITRE 2 — Section : Approche réactive au vieillissement
% ============================================================

\section{Approche réactive au vieillissement}
\label{sec:soh}

Le chapitre précédent a établi, pour chacun des trois composants de stockage, une définition scalaire de l'état de santé ainsi que les méthodes permettant de l'estimer en ligne. Cet état de santé correspond à la capacité résiduelle relative pour la batterie et à la perte (resp.\ augmentation) relative de puissance à courant constant pour la pile à combustible (resp.\ l'électrolyseur)~; il est estimé par le suivi de la résistance interne pour la batterie et par des observateurs de type filtre de Kalman étendu pour les composants hydrogène. Ces estimations sont supposées disponibles, dans la suite, à une cadence au moins hebdomadaire, ce qui est largement suffisant au regard des constantes de temps du vieillissement (plusieurs mois à plusieurs années).

Il importe de bien distinguer ces indicateurs de SoH des indices de dégradation de la section~\ref{sec:indicateurs}. Les seconds servent à \emph{évaluer a posteriori} la durabilité d'une EMS~; les premiers servent à \emph{informer le contrôle} en temps réel. C'est ce second usage qui est considéré ici. En effet, le vieillissement réduit les capacités physiques des composants~: la capacité utile de la batterie diminue, et la puissance maximale délivrable par les composants hydrogène à courant donné se dégrade. Les plages de fonctionnement nominales, définies à la mise en service, deviennent donc progressivement inadaptées~; continuer à émettre des consignes calibrées sur l'état neuf reviendrait à demander des puissances physiquement irréalisables.

L'\emph{approche réactive} adoptée dans ce chapitre consiste à exploiter le SoH estimé pour \emph{actualiser en continu les bornes physiques} (puissances maximales, capacité, bande de SoC) utilisées par le contrôle de faisabilité de l'EMS. Cette actualisation est défensive et \emph{commune à toutes les stratégies} examinées~: quelle que soit la logique de répartition, le contrôle doit rester dans l'enveloppe de fonctionnement réelle, évolutive, des composants. L'information de vieillissement n'est donc utilisée ici que pour garantir la faisabilité des commandes, sans modifier la logique de répartition. Elle se distingue en cela de l'approche \emph{anticipative} du chapitre suivant, où le SoH et la RUL sont injectés dans les règles de décision pour moduler proactivement la sollicitation des composants. L'approche réactive constitue le socle commun sur lequel cette extension anticipative viendra ensuite s'appuyer.